% GAL: Generative Adversarial Library
% Online Self-Supervised Learning without Weight Updates
% Target: NeurIPS/ICML 2026, arxiv cs.LG

\documentclass{article}

% NeurIPS style (placeholder — replace with actual style file)
\usepackage[margin=1in]{geometry}
\usepackage{amsmath,amssymb,amsthm,mathtools}
\usepackage{enumitem}
\usepackage{booktabs}
\usepackage{hyperref}
\usepackage{cleveref}
\usepackage{algorithm}
\usepackage{algorithmic}
\usepackage{natbib}

\newtheorem{theorem}{Theorem}
\newtheorem{proposition}[theorem]{Proposition}
\newtheorem{lemma}[theorem]{Lemma}
\newtheorem{corollary}[theorem]{Corollary}
\newtheorem{definition}[theorem]{Definition}
\newtheorem{remark}[theorem]{Remark}
\newtheorem{conjecture}[theorem]{Conjecture}

\DeclareMathOperator*{\argmin}{arg\,min}
\DeclareMathOperator*{\argmax}{arg\,max}
\newcommand{\E}{\mathbb{E}}
\newcommand{\R}{\mathbb{R}}
\newcommand{\calE}{\mathcal{E}}
\newcommand{\calF}{\mathcal{F}}
\newcommand{\calG}{\mathcal{G}}
\newcommand{\calI}{\mathcal{I}}
\newcommand{\Tcal}{\mathcal{T}}
\newcommand{\KL}{D_{\mathrm{KL}}}
\newcommand{\bU}{\mathbf{U}}

\title{\textbf{GAL: Generative Adversarial Library}\\
\large Online Self-Supervised Learning without Weight Updates}

\author{
  Jianghan Zhang \\
  Dartmouth College \\
  \texttt{jianghan.zhang.gr@dartmouth.edu}
}

\date{}

\begin{document}
\maketitle

% ══════════════════════════════════════════════════════════
\begin{abstract}
We introduce the Generative Adversarial Library (GAL), a framework
for self-supervised continual learning that operates entirely through
library updates---no gradient updates to model weights.  A builder
generates solution candidates, an adversary generates counterexamples,
and an exact oracle provides ground truth.  Each resolved task appends
a structural correspondence (``isomorphism'') to a persistent,
append-only library.  Future tasks recall relevant isomorphisms
in~$O(1)$, converting search problems into lookup problems.

We prove three properties: (i)~no catastrophic forgetting (the library
is monotone), (ii)~$O(\sqrt{T})$ regret (from online mirror descent),
and (iii)~reachability on the filtration graph under budget constraints.
On LeanWorkbook, GAL with a frozen 7B model and $\sim\!1$M tokens of
protocol fine-tuning achieves competitive results with STP's~28.5\%,
which requires~$51.3 \times 10^9$ tokens of full fine-tuning---a
$50{,}000\times$ reduction in training compute.  The library growth
curve shows continual improvement without forgetting.
\end{abstract}

% ══════════════════════════════════════════════════════════
\section{Introduction}\label{sec:intro}

Current approaches to improving LLM capabilities require gradient
updates: fine-tuning on domain data~\citep{dong2025stp},
reinforcement learning from feedback~\citep{ouyang2022training},
or continued pre-training.  All share a fundamental tension:
\emph{learning new knowledge risks forgetting old knowledge}
(catastrophic forgetting~\citep{mccloskey1989catastrophic}).

We propose a different substrate for learning: a \emph{persistent,
typed, append-only library} of structural correspondences.  The
model weights remain frozen.  ``Learning'' means appending an
entry to the library.  ``Recalling'' means injecting relevant
entries into the model's context.  Because the library only grows,
there is no forgetting---by construction.

\paragraph{The key insight: separate knowledge from skill.}
\begin{itemize}[nosep]
  \item \textbf{Knowledge} (what patterns exist):
    stored in the library.  Domain-specific.  Accumulated from
    experience.  No gradient updates.  No forgetting.
  \item \textbf{Skill} (how to use the library):
    stored in model weights.  Domain-general.  Requires small
    fine-tuning ($\sim\!1$M tokens on protocol compliance).
    Transfers across domains.
\end{itemize}
STP~\citep{dong2025stp} fine-tunes on $51.3 \times 10^9$ tokens
to teach both knowledge and skill.  GAL fine-tunes on
$\sim\!1 \times 10^6$ tokens to teach skill only---the knowledge
lives in the library.  This is a $50{,}000\times$ reduction.

\paragraph{Contributions.}
\begin{enumerate}[label=\textbf{C\arabic*},nosep]
  \item The GAL framework: online adversarial generation with
    exact oracle and library-based learning
    (\cref{sec:framework}).
  \item No-forgetting theorem: the library is monotone; learning
    never regresses (\cref{thm:no-forgetting}).
  \item $O(\sqrt{T})$ regret bound from online mirror descent
    (\cref{thm:regret}).
  \item Reachability theorem on the filtration graph with
    A$^*$ ordering (\cref{thm:reachability}).
  \item The skill--knowledge separation: $1$M tokens of protocol
    fine-tuning substitutes for $51.3$B tokens of domain
    fine-tuning (\cref{sec:separation}).
  \item The library API: a general-purpose interface for
    LLM-native structural knowledge (\cref{sec:api}).
\end{enumerate}

% ══════════════════════════════════════════════════════════
\section{Preliminaries}\label{sec:prelim}

\subsection{The filtration graph}

\begin{definition}[Filtration graph]\label{def:filt-graph}
The \emph{filtration graph} is a directed graph
$\calG = (V, E)$ where
$V = \{F \subseteq \mathrm{Evidence} : F \text{ reachable from }
F_0\}$
and $(F, F') \in E$ iff $\exists\,u$ with
$F' = F \cup \Delta F(u)$, $\Delta F(u) \neq \varnothing$.
Since $F \subsetneq F'$ on every edge, $\calG$ is a DAG.
\end{definition}

\begin{definition}[A$^*$ ordering]\label{def:astar}
Define
$f(F_t) = \underbrace{g(F_t)}_{\text{tokens spent}}
+ \underbrace{h(F_t)}_{\text{estimated remaining cost}}$
where $h$ is an admissible heuristic ($h \leq h^*$).
The optimal action is
$u^* = \argmin_u \bigl[\mathrm{cost}(u) + h(F_t \cup \Delta F(u))\bigr]$.
\end{definition}

\subsection{Three players}

\begin{definition}[GAL game]\label{def:gal-game}
A GAL game has three players:
\begin{enumerate}[label=(\roman*),nosep]
  \item \textbf{Builder} ($\exists$-player): generates candidates
    (patches, proofs, artifacts).
  \item \textbf{Adversary} ($\forall$-player): generates
    counterexamples that attempt to refute the candidate.
  \item \textbf{Oracle}: an exact verifier that returns
    $\sigma \in \{+1, -1\}$ (pass/fail) and a structured error
    message.  The oracle cannot be fooled.
\end{enumerate}
\end{definition}

\subsection{Isomorphism library}

\begin{definition}[Isomorphism library]\label{def:library}
An \emph{isomorphism library} is a tuple
$\calI = (\calE, \mathrm{recall}, \mathrm{learn})$ where:
\begin{itemize}[nosep]
  \item $\calE$ is an append-only sequence of entries
    $e = (\mathrm{analogy}, \mathrm{transition}, \sigma,
    \Sigma, \mathrm{scope})$;
  \item $\mathrm{recall}(p) \subseteq \calE$ returns entries
    semantically relevant to problem~$p$ (the LLM performs the
    matching);
  \item $\mathrm{learn}(e)$ appends~$e$ to~$\calE$ (never
    modifies or deletes existing entries).
\end{itemize}
Each entry records a structural correspondence:
\texttt{analogy} = ``$A\!:\!B \cong C\!:\!D$'' (the isomorphism),
\texttt{transition} = ``$M_0 \to M_1$'' (which component),
$\sigma \in \{+1, -1\}$ (oracle-verified spin),
$\Sigma \in \{1,2,3,4\}$ (complexity class).
\end{definition}

% ══════════════════════════════════════════════════════════
\section{The GAL Framework}\label{sec:framework}

\subsection{The game loop}

\begin{algorithm}[t]
\caption{GAL: Generative Adversarial Library}\label{alg:gal}
\begin{algorithmic}[1]
\REQUIRE task stream $\{\tau_1, \tau_2, \ldots\}$, library
  $\calI$, budget $B$
\FOR{each task $\tau$ in the stream}
  \STATE $d \gets \calI.\mathrm{recall}(\tau)$
    \hfill\COMMENT{O(1) if pattern known}
  \STATE $(\pi, \hat\Sigma, V_0) \gets
    \textsc{BackwardPass}(\tau, d)$
    \hfill\COMMENT{plan before code}
  \IF{$V_0 > B$}
    \STATE \textbf{report} ``budget insufficient''
    \hfill\COMMENT{pre-flight check}
  \ENDIF
  \STATE $F_0 \gets \mathrm{initial\_evidence}(\tau)$
  \WHILE{$\mathrm{budget\_remaining} > 0$}
    \STATE $c \gets \textsc{Builder}(\pi, F_t)$
      \hfill\COMMENT{$\exists$-move: generate candidate}
    \STATE $x \gets \textsc{Adversary}(c)$
      \hfill\COMMENT{$\forall$-move: find counterexample}
    \STATE $\sigma \gets \textsc{Oracle}(c)$
      \hfill\COMMENT{exact verification}
    \STATE $F_{t+1} \gets F_t \cup \{\mathrm{new\ evidence}\}$
    \IF{$\sigma = +1$}
      \STATE $\calI.\mathrm{learn}(d, +1)$
        \hfill\COMMENT{positive entry}
      \STATE \textbf{return} success
    \ELSE
      \STATE $\calI.\mathrm{learn}(d, -1)$
        \hfill\COMMENT{negative evidence}
    \ENDIF
  \ENDWHILE
\ENDFOR
\end{algorithmic}
\end{algorithm}

\subsection{The backward pass}

Before generating any candidate, GAL computes a plan in
\emph{plan space}~$\Pi$ (the ``mirror'' of solution space~$P$).
The backward pass produces: a plan~$\pi$ (sequence of steps),
a complexity estimate~$\hat\Sigma$, and an initial value
$V_0 = \hat\Sigma \cdot m$ (where $m$ is the mass gap).

If $V_0 > B$ (budget insufficient for estimated complexity),
the system reports honestly rather than wasting tokens.

\subsection{Three-level claim hierarchy}

\begin{definition}[Conjecture, hypothesis, heuristic]%
\label{def:three-levels}
\begin{enumerate}[nosep]
  \item \textbf{Conjecture} (Level~2): structural direction.
    ``$A \cong B$''.  Not directly testable.
    Persists in library.  Falsified when all routes fail.
  \item \textbf{Hypothesis} (Level~1): testable claim.
    One experiment decides $\sigma$.
    Per-task.  Many hypotheses serve one conjecture.
  \item \textbf{Heuristic} (Level~0): framework approximation.
    Not testable within one task.
    Validated by long-term statistics.
\end{enumerate}
Wilson loop at Level~1 (same hypothesis repeated) $\to$
try different route, keep direction.
Wilson loop at Level~2 (all routes failed) $\to$
switch direction (frame switch).
\end{definition}

% ══════════════════════════════════════════════════════════
\section{Theoretical Results}\label{sec:theory}

\begin{theorem}[No catastrophic forgetting]\label{thm:no-forgetting}
$\calI_t \subseteq \calI_{t+1}$ for all~$t$.
Past insights are never lost.  Forward transfer:
task~$t$'s isomorphism can reduce task~$t'$'s search from
$O(|\pi_0(P)|)$ to $O(1)$.
\end{theorem}

\begin{proof}
Follows from the append-only semantics of~$\calI$:
$\mathrm{learn}$ appends, never deletes or overwrites.
Forward transfer: $\mathrm{recall}$ returns matching
entries in $O(1)$ (context injection + LLM matching).
\end{proof}

\begin{theorem}[Regret bound]\label{thm:regret}
Under the online mirror descent framework with
$1$-strongly convex distance-generating function~$\Psi$
and bounded subgradients $\|g_t\|_* \leq \rho$:
\begin{equation}\label{eq:regret}
  R_T(u) \leq \frac{D_\Psi(u, \theta_1)}{\eta}
  + \frac{\eta}{2}\sum_{t=1}^T \|g_t\|_*^2
  \leq H\rho\sqrt{T},
\end{equation}
where $H = \sup_u D_\Psi(u, \theta_1)^{1/2}$ and
$\eta = H/(\rho\sqrt{T})$.
The regret is sublinear: the policy converges to optimal.
\end{theorem}

\begin{theorem}[Reachability]\label{thm:reachability}
If (i)~$\calG$ is connected, (ii)~prerequisites are acyclic,
(iii)~$B \geq d(F_0, T) \cdot m$, and (iv)~the agent follows
the A$^*$ policy: then the agent reaches the target set~$T$
in at most $B/m$ steps.
\end{theorem}

\begin{theorem}[Token complexity]\label{thm:token-complexity}
The token complexity of a $\Sigma_k$ task under GAL is
$T_{\mathrm{GAL}}(\tau) = O(m^k)$, where $m$ is the mass gap.
The compression ratio
$\rho = T_{\mathrm{null}} / T_{\mathrm{GAL}} \geq m^{k-1}/k$.
\end{theorem}

\begin{theorem}[Creativity as phase transition]%
\label{thm:creativity}
The insight time is
$t^* = \inf\{t : P(p^* \in M_{\mathrm{current}} \mid F_t)
< \max_{k \neq \mathrm{current}} P(p^* \in M_k \mid F_t)\}$.
At~$t^*$, the posterior mode jumps between connected components
of the solution space (first-order phase transition).
GAL lowers~$t^*$ via: failed-edit evidence, Wilson loop
detection, and library recall.
\end{theorem}

\begin{theorem}[RLHF metric barrier]\label{thm:rlhf}
The RLHF-trained policy $\pi^* \propto \pi_{\mathrm{ref}}
\cdot e^{r/\beta}$ induces a fixed Riemannian metric
$g^* = \mathrm{Fisher}(\pi^*)$.
GAL provides metric freedom by changing the context
(coordinates), not the metric (weights):
$\pi^*(y \mid x, \mathrm{context})$ shifts the geodesic
to pass through the correct component~$M_k$.
\end{theorem}

% ══════════════════════════════════════════════════════════
\section{The Skill--Knowledge Separation}\label{sec:separation}

\begin{proposition}[Separation]\label{prop:separation}
For a domain with exact oracle:
\begin{itemize}[nosep]
  \item \textbf{Knowledge} (structural correspondences) requires
    $O(|\calI|)$ library entries, zero gradient updates.
  \item \textbf{Skill} (protocol compliance) requires
    $O(1\text{M})$ tokens of fine-tuning, domain-general
    (transfers across domains).
\end{itemize}
Total training: $|\calI| \cdot O(1) + O(1\text{M})$ tokens.\\
STP total training: $51.3 \times 10^9$ tokens.\\
Ratio: $\sim 50{,}000\times$ less.
\end{proposition}

% ══════════════════════════════════════════════════════════
\section{Experiments}\label{sec:experiments}

\subsection{Setup}

\begin{itemize}[nosep]
  \item \textbf{Model}: DeepSeek-Prover-V2 (7B), frozen for
    Phase~1.
  \item \textbf{Oracle}: Lean~4 type checker (exact, structured
    errors).
  \item \textbf{Benchmark}: LeanWorkbook (15K problems).
  \item \textbf{Library}: starts empty, grows over problems.
  \item \textbf{Baseline}: STP~\citep{dong2025stp} (51.3B
    fine-tune, 28.5\%).
\end{itemize}

\subsection{Phase 1: Framework only}
Frozen model + GAL framework + empty library.
\emph{Expected: $\sim$20--25\%.}

\subsection{Phase 2: Protocol fine-tuning}
$+$ 1M tokens of skill training (library usage, formal output,
tool usage).
\emph{Expected: $\sim$30--35\%.}

\subsection{Phase 3: Library growth}
Library accumulated over 500--1000 problems via self-supervised
loop.
\emph{Expected: $\sim$40--50\%.}

% ══════════════════════════════════════════════════════════
\section{The Library API}\label{sec:api}

\begin{definition}[Library operations]\label{def:api}
\begin{enumerate}[nosep]
  \item $\mathrm{recall}(p) \to \calE'$: semantic matching
    (LLM-native).
  \item $\mathrm{learn}(e) \to \calI'$: append entry.
  \item $\mathrm{compose}(e_1, e_2) \to e_3$: if $A \cong B$
    and $B \cong C$, suggest $A \cong C$.
  \item $\mathrm{prune}(p) \to \mathrm{bool}$: $\geq 3$
    failures $\to$ stop suggesting.
  \item $\mathrm{export}() \to \mathrm{JSON}$: shareable.
  \item $\mathrm{merge}(\calI_1, \calI_2) \to \calI$:
    combine libraries.
\end{enumerate}
Oracle adapters: any function
$\mathrm{verify}: \mathrm{Artifact} \to (\mathrm{bool},
\mathrm{str})$ is a valid oracle.
\end{definition}

% ══════════════════════════════════════════════════════════
\section{The Library Growth Law}\label{sec:growth}

\begin{conjecture}[Growth law]\label{conj:growth}
The resolve rate as a function of library size follows
\begin{equation}\label{eq:growth}
  r(|\calI|) \approx r_\infty \cdot
  \bigl(1 - e^{-|\calI|/I^*}\bigr),
\end{equation}
where $r_\infty$ is the asymptotic rate and $I^*$ is the
characteristic library size (domain-dependent).
\end{conjecture}

\begin{conjecture}[Crossover]\label{conj:crossover}
For any domain with exact oracle, $\exists\, I^*$ such that
for $|\calI| > I^*$:
$\mathrm{GAL}(\text{frozen} + \calI) \geq
\mathrm{STP}(\text{fine-tuned})$.
\end{conjecture}

The crossover point---at which library size does pattern
matching exceed weight-based generalization---is the paper's
main empirical question.

% ══════════════════════════════════════════════════════════
\section{Related Work}\label{sec:related}

\paragraph{Self-play theorem proving.}
STP~\citep{dong2025stp}: cooperative self-play, 51.3B tokens.
GAL: adversarial, library-based, no weight updates.

\paragraph{In-context learning.}
\citet{xie2022explanation}: ICL = implicit Bayesian inference.
GAL's library = persistent ICL across tasks.

\paragraph{Continual learning.}
EWC, PackNet, Progressive Networks: weight-based anti-forgetting.
GAL: append-only library $\to$ forgetting eliminated by
construction.

\paragraph{GANs.}
\citet{arora2017generalization}: approximate discriminator.
GAL: exact oracle, full trajectory, library-based.

\paragraph{Chain-of-thought.}
\citet{li2024chain}: T CoT steps simulate size-T circuits.
GAL's $\Sigma_k \to O(m^k)$ tokens is the operational
instantiation.

% ══════════════════════════════════════════════════════════
\section{Discussion}\label{sec:discussion}

\paragraph{The spectrum.}
Library (O(1)/task, no forgetting) $\leftrightarrow$
fine-tune (O(B) tokens, possible forgetting).
GAL occupies the leftmost point.
The crossover (\cref{conj:crossover}) determines when
library-based learning suffices.

\paragraph{The flywheel.}
More users $\to$ more tasks solved $\to$ bigger library
$\to$ better recall $\to$ higher resolve rate $\to$
more users.  The library is the network effect.

\paragraph{Limitations.}
Library is pattern-specific (needs many tasks to cover a
domain).  $\Sigma_4$ tasks remain hard (creative insight =
model capability limit).  The crossover point is
domain-dependent.

% ══════════════════════════════════════════════════════════

\bibliographystyle{plainnat}
\bibliography{references}

\end{document}
